\section{Uniform rates and the order of the mixing depth}
\label{sec:quantitative-families}

The dichotomy supplies generating sets at bounded depth. Diameter estimates
then control their contraction uniformly over all valid templates.
In particular, the unobstructed exterior
classes of $\PGL_2(q)$ have base depth $2$ or $3$, uniformly in $q$;
the issue here is the strength of contraction as the field grows.
All logarithms in this section are natural
unless their base is displayed. For a tree $\tau$, the law $\mu_\tau$
is unconditional.

\begin{theorem}[Diameter-to-mixing transfer]
\label{thm:family-rate}
Consider a family of pairs $(G,\mathcal C)$ satisfying case~\textup{(2)} of
\Cref{thm:dichotomy}, with exact levels $D\le D_0$. Suppose there are
constants $K,a>0$, independent of the pair, such that every symmetric
generating set $S$ of $G$ satisfies
$\mathrm{diam}(G,S)\le K(\log\abs G)^a$.
Put $c=(8\cdot16^{D_0}K^2)^{-1}$. Every valid configurable tree $\tau$ of
level $m\ge D$ satisfies
\begin{equation}\label{eq:family-rate}
 \dtv(\mu_\tau,U_G)\le\frac{\sqrt{\abs G-1}}2
       \exp\left(-\frac{c4^m}{(\log\abs G)^{2a}}\right).
\end{equation}
For $D_T$, add the conditioning probability $p$ from \Cref{lem:conditioning}.
The leaves need not be involutions.
\end{theorem}
\begin{proof}
Every sub-block of a valid tree is valid. Its level-$D$ sub-blocks therefore
generate $G$ by \Cref{thm:dichotomy}. Each has $N=4^D$ expanded occurrences,
including inverse occurrences. Fix a non-trivial irreducible $\rho$.
For each level-$D$ sub-block $\tau_0$, \Cref{prop:stability} gives
\[
 \opn{M_{\tau_0}(\rho)}\le1-\frac1{8\cdot4^D K^2(\log\abs G)^{2a}}.
\]
Expanding the recursion down to this level produces $4^{m-D}$ factors,
counting adjoints, all with that norm bound. Submultiplicativity and
$1-x\le e^{-x}$ yield
\[
 \opn{M_\tau(\rho)}\le
 \exp\left(-\frac{4^m}{8\cdot16^D K^2(\log\abs G)^{2a}}\right).
\]
Apply Diaconis--Shahshahani as in \Cref{prop:quant}, and use $D\le D_0$.
\end{proof}

For $0<\eta<1/2$, a sufficient condition for error at most $\eta$ is
$m\ge D$ and
\begin{equation}\label{eq:family-accuracy}
 4^m\ge c^{-1}(\log\abs G)^{2a}
       \left(\tfrac12\log(\abs G-1)+\log\frac1{2\eta}\right).
\end{equation}
Thus fixed error, or error $\abs G^{-b}$ for fixed $b>0$, requires only
$O(\log\log\abs G)$ depth and $O((\log\abs G)^{2a+1})$ expanded length.

\begin{proposition}[Support lower bound and optimal order]
\label{prop:depth-lower}
For any level-$m$ independent-read tree $\tau$,
\[
 \dtv(\mu_\tau,U_G)\ge\max\{0,1-2^{4^m}/\abs G\}.
\]
Consequently, for any fixed $0<\eta<1/2$, in an unbounded family satisfying
\Cref{thm:family-rate}, the least depth at which every valid tree has
error at most $\eta$ is $\Theta(\log\log\abs G)$, uniformly over the family.
\end{proposition}
\begin{proof}
There are $4^m$ independent input bits, hence at most $2^{4^m}$ possible
outputs. The complement of their support has uniform mass at least
$1-2^{4^m}/\abs G$ and output mass zero. Error at most $\eta$ therefore
requires $4^m\ge\log_2((1-\eta)\abs G)$.
Valid trees exist at every level by \Cref{lem:central}, since $G$ is
non-solvable. This rules out vacuous small levels. The upper bound is
\eqref{eq:family-accuracy}.
\end{proof}

The upper and lower bounds match in their order of depth, up to constant factors.
On the obstructed side, the subgroup $H<G$ in \Cref{thm:dichotomy}(1)
supplies at every level a valid tree $\tau$ supported in $H$, and hence
$\dtv(\mu_\tau,U_G)\ge1-\abs H/\abs G\ge1/2$.

\begin{corollary}[Rank one and minimal simple groups]
\label{cor:family-examples}
The uniform rate \eqref{eq:family-rate} and the optimal order of depth in
\Cref{prop:depth-lower} apply to:
\begin{enumerate}[nosep,leftmargin=1.6em]
\item the unobstructed involution classes of $\PSL_2(q)$ and exterior
classes of $\PGL_2(q)$ in \Cref{thm:rankone-classification};
\item all finite minimal simple groups with any non-trivial conjugacy
class, uniformly over both the group and the class.
\end{enumerate}
\end{corollary}
\begin{proof}
Pyber--Szab\'o's bounded-rank diameter theorem~\cite[Theorem 2 of the
preprint]{PyberSzabo2016} supplies constants $K,a$ for every symmetric
generating set of a simple group of Lie type of bounded rank.
For $G=\PGL_2(q)$, take a symmetric generating set $S$ and $t\in S$
outside $L=\PSL_2(q)$. With transversal $\{1,t\}$, Schreier generators
of $L$ have length at most $3$ in $S$; their inverses satisfy the same
bound. Writing this symmetric set as $S_L$ gives
\[
 \mathrm{diam}(G,S)\le3\,\mathrm{diam}(L,S_L)+1.
\]
Thus the same type of diameter bound holds for $G$, with adjusted absolute
constants. The rank-one levels are at most $3$.
For (2), Thompson's list in \Cref{cor:minsimple} has bounded Lie rank
(at most $2$, including the underlying type of the Suzuki groups), and
\Cref{cor:minsimple} gives $D\le5$. Apply \Cref{thm:family-rate} with
$D_0=3$ or $5$, respectively, and then \Cref{prop:depth-lower}.
\end{proof}

\subsection{Uniform expansion outside exceptional primes}

Over prime fields, strong uniform expansion gives constant gaps for every
generating block of bounded depth. The prime-field input goes back to
Breuillard--Gamburd~\cite[Corollary 1.1]{BreuillardGamburd2010}.
Becker--Breuillard~\cite[Theorem 1.1 and Corollary 3.2]{BeckerBreuillard2026}
extend uniform expansion to bounded-rank families with bounded field degree,
outside exceptional sets of characteristics.

\begin{corollary}[Uniform gaps for almost all primes]
\label{cor:almost-all-primes}
There is a set $\mathcal E$ of relative density zero among the primes
such that the following hold for every prime $p\ge5$ outside $\mathcal E$.
\begin{enumerate}[nosep,leftmargin=1.6em]
\item For each fixed $K\ge1$ there is $a_K>0$, independent of $p$, such that
every involutive block $\tau$ of level $k\le K$ whose leaves generate
$G=\PSL_2(p)$ or $\PGL_2(p)$ satisfies
$\opn{M_\tau(\rho)}\le1-a_K$ for every non-trivial irreducible $\rho$.
\item Let $\mathcal C$ be an unobstructed involution class from
\Cref{thm:rankone-classification}: the exterior class in $\PGL_2(p)$,
or the involution class in an unobstructed $\PSL_2(p)$.
There is an absolute $c>0$ such that every valid tree $\tau$ of level $m\ge D$
over $\mathcal C$ satisfies
\begin{equation}\label{eq:almost-all-prime-rate}
 \opn{M_\tau(\rho)}\le e^{-c4^m},\qquad
 \dtv(\mu_\tau,U_G)\le\frac{\sqrt{\abs G-1}}2 e^{-c4^m}.
\end{equation}
The operator bound holds for every non-trivial irreducible $\rho$.
\item If $p\equiv\pm3\pmod8$, every pair of exterior subgroups
$A,B\cong S_4$ generating $\PGL_2(p)$ satisfies
$\opn{P_A^\rho P_B^\rho}\le1-b$ for every non-trivial irreducible $\rho$,
with an absolute $b>0$ independent of the pair and of $p$.
\end{enumerate}
\end{corollary}
\begin{proof}
Strong uniform expansion supplies $\varepsilon\in(0,1]$ and
$\mathcal E$ such that $|SX|\ge(1+\varepsilon)|X|$ whenever
$|X|\le|G|/2$ and $S$ is a symmetric generating set containing $e$,
for both groups $G$. Starting with
\cite[Corollary 1.1]{BreuillardGamburd2010}, pass from $\SL_2(p)$ to
its quotient $\PSL_2(p)$ and then to its index-two overgroup
$\PGL_2(p)$; uniform expansion is preserved under these operations
\cite[Lemma 3.9]{BeckerBreuillard2026}. Both sources use vertex
expansion: the proof of \cite[Corollary 1.1]{BreuillardGamburd2010}
produces a single $s\in S$ with $|sX\setminus X|\ge c|X|$, and
\cite[(1.1)]{BeckerBreuillard2026} states the displayed form directly.

For (1), fix such a block $\tau$ and list its $r=2^k$ leaf positions with their multiplicities.
Their lazy walk assigns each distinct non-identity leaf probability at
least $1/(2r)$. The expansion inequality gives at least
$\varepsilon|X|$ outgoing edges from any $X$ with $|X|\le|G|/2$,
so its conductance is at least $\varepsilon/(2r)$.
Cheeger's inequality gives spectral gap
$\delta\ge\varepsilon^2/(8r^2)$, including these multiplicities.
Each leaf occurs $2^k$ times in the expanded involutive block, so this
is also its expanded-leaf walk. Put $N=4^k=r^2$. By
\Cref{prop:cayley-comparison},
\[
 1-\opn{M_\tau(\rho)}^2
 \ge\frac{\varepsilon^2}{4N(N+1)}.
\]
Thus one may take $a_K=\varepsilon^2/(8\cdot4^K(4^K+1))$.

For (2), $D\le3$, and every level-$D$ sub-block of a valid tree $\tau$
generates $G$ by \Cref{thm:dichotomy}. Applying (1) at this fixed
base level and expanding the recursion gives
\[
 \opn{M_\tau(\rho)}\le(1-a_3)^{4^{m-D}}
 \le e^{-(a_3/64)4^m}.
\]
Take $c=a_3/64$ and apply the Diaconis--Shahshahani bound.

For (3), fix a non-trivial irreducible $\rho$ on $V$, and write
$P_A=P_A^\rho$, $P_B=P_B^\rho$.
The walk $\nu=(U_A+U_B)/2$ has generating support
$A\cup B$ and assigns each support element probability at least $1/48$.
The same conductance argument gives spectral gap at least
$\eta=\varepsilon^2/(2\cdot48^2)$. Its Fourier operator is
$(P_A+P_B)/2$, which is positive semidefinite.
For a unit vector $v\in V^A$,
$1+\norm{P_Bv}^2\le2(1-\eta)$.
Hence $\opn{P_AP_B}\le\sqrt{1-2\eta}\le1-\eta$;
if $V^A=0$ the norm is zero. Take $b=\eta$.
\end{proof}

The all-pairs octahedral programme was formulated in the July 2026
preprint~\cite[Conjectures 6.9 and 7.8, Remark 7.9]{ZermenoRankOneZenodo},
which posed a sharp angular bound and a uniform non-concentration target.
\Cref{cor:almost-all-primes} gives qualitative gaps for all generating
exterior pairs and all ambient-generating level-$2$ involutive blocks
outside $\mathcal E$.

\subsection{Few exceptional configurations over arbitrary fields}

\begin{corollary}[Exceptional generating blocks at fixed depth]
\label{cor:few-exceptional-blocks}
Fix integers $R,k\ge1$ and $\delta>0$.
There is $a=a(R,k,\delta)>0$ such that, for every finite simple group
$G(q)$ of Lie type of rank at most $R$, all but at most
$O_{R,k,\delta}(q^\delta)$ simultaneous $G(q)$-conjugacy classes of
generating involutive level-$k$ configurations $\tau$ satisfy
\[
 \opn{M_\tau(\rho)}\le1-a
 \qquad\text{for every non-trivial irreducible }\rho.
\]
The same conclusion holds for $\PGL_2(q)$, $q\ge4$.
Restricting to valid configurations, or to specified conjugacy classes
of involution leaves, preserves this absolute bound on exceptions.
\end{corollary}
\begin{proof}
The dimension-zero theorem of Becker--Breuillard
\cite[Theorems 1.2 and 4.1]{BeckerBreuillard2026}, with $2^k$ entries,
gives an expansion constant $\varepsilon=\varepsilon(R,k,\delta)>0$
outside the stated number of simultaneous conjugacy classes of
generating tuples. The finite simple groups have the required
bounded-dimensional algebraic embeddings; $\PGL_2(q)$ is covered
by Theorem~4.1 in the adjoint group of type $A_1$, since it contains
$\PSL_2(q)$ and therefore satisfies the density hypothesis of that
theorem for $q$ large.
Finitely many small fields are absorbed into the implied constant.
For an involutive tuple put $N=4^k$. Its lazy leaf walk, with
multiplicities, has gap at least $\varepsilon^2/(8N)$ by the
conductance argument in \Cref{cor:almost-all-primes}. By
\Cref{prop:cayley-comparison} one may take
$a=\varepsilon^2/(8N(N+1))$, after reducing $\varepsilon$ to at most $1$.
Conjugating every leaf conjugates the block operator, so the transfer
respects exactly the equivalence relation counted by the source theorem.
\end{proof}

Here the depth is fixed and the field degree is unrestricted. The
exception count is absolute; a relative-density statement in a
restricted family additionally requires a lower bound on its size.
The next section obtains the sharp natural-pair constant for every
prime in its stated congruence classes.
